\documentclass[12pt]{article}
\usepackage[T1]{fontenc}
\usepackage[utf8]{inputenc}
\usepackage{lmodern}
\usepackage{textcomp}
\usepackage{enumitem}
\usepackage{geometry}
\usepackage{graphicx}
\usepackage{amsmath, amssymb}
\geometry{margin=1in}
\begin{document}
\begin{center}
Economics - Graduate Level Assessment 14 \
Total Marks: 40 \
Answer all questions.
\end{center}

\section*{Multiple Choice (10 marks)}
\begin{enumerate}[label=\arabic*.]
\item Scarcity is best defined as
\begin{enumerate}[label=(\alph*)]
    \item the opportunity cost of not pursuing a given course of action.
    \item the difference between unlimited wants and limited economic resources.
    \item the difference between unlimited resources and unlimited economic wants.
    \item the opportunity cost of pursuing a given course of action.
    \item the difference between the total benefit of an action and the total cost of that action.
\end{enumerate}
\item Suppose real GDP increases. We can conclude without a doubt that
\begin{enumerate}[label=(\alph*)]
    \item the average wage is higher.
    \item the national debt is lower.
    \item production is higher.
    \item prices are higher.
    \item employment is higher.
\end{enumerate}
\item Labor productivity and economic growth increase if
\begin{enumerate}[label=(\alph*)]
    \item a nation discourages investment in technology.
    \item a nation decreases spending on public services.
    \item a nation ignores societal barriers like discrimination.
    \item a nation promotes monopolies in industry.
    \item a nation imposes tariffs and quotas on imported goods.
\end{enumerate}
\item The fixed effects panel model is also sometimes known as
\begin{enumerate}[label=(\alph*)]
    \item The pooled cross-sectional model
    \item The random effects model
    \item The fixed variable regression model
    \item The instrumental variables regression model
    \item The generalized least squares approach
\end{enumerate}
\item If the standard of living increases we can conclude that
\begin{enumerate}[label=(\alph*)]
    \item output must have increased proportionally more than population.
    \item the cost of living and population must have increased.
    \item the cost of living must have decreased.
    \item population must have decreased.
    \item output and population must have increased.
\end{enumerate}
\end{enumerate}

\section*{True / False (10 marks)}
\textit{Answer True or False}
\begin{enumerate}[label=\arabic*.]
\setcounter{enumi}{5}
\item True or False: A natural monopoly achieves allocative efficiency when price is equal to marginal cost.
\item True or False: The slope of the consumption function is equal to the marginal propensity to consume (MPC).
\item True or False: The bank's actual reserves can be calculated as \$50,000 based solely on the required reserve ratio of 20\%.
\item True or False: Shocks in a stationary autoregressive process will lead to a permanent change in the trend of the data.
\item True or False: A decrease in the reserve ratio will result in an increase in excess reserves held by banks.
\end{enumerate}

\section*{Long Form Response (20 marks)}
\textit{Answer each question with a clear and complete explanation.}
\begin{enumerate}[label=\arabic*.]
\setcounter{enumi}{10}
\item Discuss the effects of contractionary fiscal policy in the United States on domestic interest rates and the international value of the dollar, providing a detailed analysis of the mechanisms involved.
\item Discuss the characteristics of pure public goods, focusing on the implications of zero marginal cost for additional consumption by citizens.
\end{enumerate}

\vfill
\noindent\textit{End of Paper}
\end{document}